\cotefile{78}
\scanpage{103}
b) Tout sommet $s=\varphi_S(\underline{s})$ appartient à
exactement deux arêtes. Soient en effet
$\underline a,\underline a'\in\underline A$ les arêtes incidentes à $\underline s$, $a,a'$
leurs images dans $A$, \deleted{\unsure{on}} sait que $a\ne a'$,
$s\in a$, $s\in a'$. Montrons que si $a_0\in\varphi_A(\underline a_0)$
\deleted{est} contient $s$, alors $a_0=a$ ou $a_0=a'$,
\deleted{\unsure{Prenons}} Soient $\deleted{\underline s_0}\underline{\unsure{t|s}}_0,\underline t_1$ les sommets incidents
à $\underline a_0$, \deleted{donc} $t_0,t_1\in S$ leurs images, donc
\deleted{\unsure{$s,t,t'$ sont|$s_0\ne s_1$, et}} \inserted{$t_0\ne t_1$, et} $a_0=\text{\unsure{dr|Dr}}\,(t_0,t_1)\ni s$, et par a) on
doit avoir \deleted{\unsure{$t\in s$}} $s=t_0$ ou $s=t_1$, \deleted{\unsure{et par|disons}}
\deleted{$s$ est inci} $s=\underline t_0$. Soit $\underline r=(\underline t_0,\underline a_0)$,
on a $\underline s=\underline u^{n}\underline t_0$, pour $n\in\mathbb Z$ convenable,
donc $u^nt_0=\deleted{s}=s_0$ \deleted{\unsure{d'où}} $u^n=\mathrm{id}$, donc
\deleted{\illegible{1}} $u^n(a_0)\deleted{\unsure{=ua}} \inserted{=a_0}$. Or $\underline{s}_0$ étant incident
à $\underline a_0$, $\underline s=\underline u^n(\underline s_0)$ l'est à $\underline u^n(\underline a_0)$, donc
$\underline u^n(\underline a_0)$ \deleted{\unsure{$\underline a$}} est égal à $\underline a$ ou $\underline a'$, donc
$u^n(a_0)$ \inserted{$=a_0$} est égal à $a$ ou $a'$, \deleted{donc} cqfd.

c) \deleted{\unsure{les}} $S$ engendre $E$ : connu.

d) $G=\mathrm{Aut}(E,S,A)$ est transitif sur
l'en\inserted{$R$} des repères de $(S,A)$. Dans
la dém. de a) \inserted{(ou comme cas immédiat de a))} on a vu que $R$ est
% archive sheet number: 102
\scanpage{104}
l'image canonique de $\underline R$. D'ailleurs
$\underline G$ opérant sur $\underline R$ transitivement,
il opère transitivement sur $R$, donc
son image $G$ dans $\mathrm{Aff}(E,S,A)$ opère transitivement
sur $R$. On trouve en $=$
temps que (comme $\mathrm{Aut}(E,S,A)$ opère de
façon simplement transitive) que
\[
G\overset{\text{\unsure{dfn|défn|déf}}}{=}\mathrm{Im}\,\underline G=\mathrm{Aut}(E,S,A)
\]

e) Pour prouver que $\underline\Pi\to\Pi_{\deleted{\unsure{a|0}}}$ est
un hom, il faut voir que $\underline R\unsure{-|\to}R$
est compatible avec $\sigma_0,\sigma_1$ -- on l'a
vu dans la dém. de a) et \uline{b}).
% archive sheet number: 103
\scanpage{105}
\begin{center}\uline{Fidélité}\end{center}

% slanted
\margin{(ancienne version)}
\deleted{\uline{\unsure{Corollaire}}} \inserted{Proposition} Soit une réalisation géométrique
régulière $\varphi$ \inserted{(non dégénérée (et propre))} d'un polygone combinatoire
$\Pi\inserted{=(\underline S,\underline A,\underline R)}$ dans un plan affine $E$ \inserted{sur un corps $k$}. Conditions équiv.:

a) $\varphi$ est fidèle

b) $\varphi_S:\underline S\to E$ injective

b$'$) $\varphi_A:\underline A\to\mathrm{Dr}(E)$ injective

c) $\mathrm{Aut}(\Pi)=G\unsure{-|\to}\mathrm{Aff}(E)$ injective

d) \deleted{[Si $\Pi$ orienté, d'où $\underline u\in\mathrm{Aut}(\Pi)$, et $u\in\mathrm{Aff}(E)$]}
\deleted{$\mathrm{Aut}^{+}(\Pi)$} $G^+\unsure{-|\to}\mathrm{Aff}(E)$ injective

\uline{dém} a) $\Rightarrow$ b) \& b$'$) trivial (\unsure{compte|cpte} tenu que
$\forall\underline s,\underline s'\in\underline S$, $\underline s\ne\underline s'$, $\exists\underline a\in\underline A$ qui est incident à $\underline s$ et
non à $\underline s'$ -- car $\mathrm{card}\,S=\mathrm{card}\,A\geq3$ ici ! -- \inserted{et dualement.}

b) $\Rightarrow$ c), \inserted{b$'$) $\Rightarrow$ c)} \unsure{triviales|trivial|triviaux}, puisque $G\to\unsure{\mathfrak S|\mathcal S}_S$, \inserted{$G\to\mathfrak S_A$} injectives.

\deleted{\unsure{Prouvons}} c) $\Rightarrow$ d) trivial. Prouvons
d) $\Rightarrow$ b), i.e. non b $\Rightarrow$ non d). Choisissons
\deleted{un gén.} repère $r$ de $\Pi$, d'où un gén.
\deleted{\unsure{u de}} $\underline u=\underline\sigma_0\underline\sigma_1$ de $\underline G=\mathrm{Aut}(\Pi)$, d'où $u\in\mathrm{Aff}(E)$.
\deleted{On aura} Soit $n=\mathrm{Card}\,\underline S$ (\deleted{de sorte que \unsure{les}} \inserted{$3\leq n\leq+\infty$})
\deleted{Dire que} \inserted{Alors} non b signifie $\exists d<n$ \deleted{tel que}
\deleted{$\underline s_d$} et $i\in\mathbb Z$ tels que $s_i=s_{i+d}$ i.e.
$u^is_0=u^{i+d}s_0$ i.e. \deleted{\unsure{$s$}} $u^ds_0=s_0$ ce qui
implique, en multipliant par $u^j$ et
notant $u^ju^ds_0\unsure{-|=}u^du^js_0=u^ds_j$, que
$u^ds_j=s_j$ $\forall j\in\mathbb Z$, et comme les $s_j$ engendrent
$E$, on aura $u^{d}=\mathrm{id}$, \deleted{\unsure{donc}} or $\underline u^d\ne1$, donc
non d). On a ainsi prouvé
% archive sheet number: 104
\scanpage{106}
% A vertical pen stroke groups the statement at the left.
\uline{Corollaire} Avec les notations précédentes, si
$\underline v\in G^+$, $s\in S=\mathrm{Im}(\underline S\to E)$, alors
\[
\underline v.s=s\Longleftrightarrow v=\mathrm{id}\qquad (v=\varphi_G(\underline v)\in\mathrm{Aff}(E))
\]

\deleted{[\unsure{Mais} b) signifie que $\underline v s_{\unsure{0}}=s_{\unsure{0}}$ implique $\underline v\underline s_0=\underline s_0$]}
Ainsi,
\diagram{\begin{tikzcd}[ampersand replacement=\&,column sep=small,row sep=tiny]
\text{a)}\arrow[dr,Rightarrow]\&{}\&{}\&{}\&{}\\
{}\&\text{b$'$)}\arrow[dr,Rightarrow]\&{}\&{}\&{}\\
{}\&{}\&\text{b)}\arrow[r,Leftrightarrow]\&\text{c)}\arrow[r,Leftrightarrow]\&\text{d)}
\end{tikzcd}}
\deleted{d'ailleurs}

\deleted{et b) $\unsure{-|\Rightarrow}$ b$'$), d'ailleurs \unsure{$\neg$ b)} signifie \deleted{\unsure{implique}}.
$\exists d<n$ tel que \deleted{$u^ds_0=s_0$} \inserted{$u^d=\mathrm{id}$} implique que
$u^da_0=a_0$ i.e. $a_d=a_0$ donc non b$'$), donc}

Il reste à prouver b), c), d) impliquent a).
Soit donc $\underline s\in\underline S$, $\underline a\in\underline A$ supposons $s\in a$
prouvons $\underline s$ incident à $\underline a$. Soient $\underline s_0,\underline s_1$ les
deux sommets \deleted{\unsure{de}} incidents à $\underline a$ \inserted{(distincts)}, donc $s_0,s_1\in a$,
\inserted{et $s_0\ne s_1$,} il faut prouver que \deleted{\unsure{si $s$}} $s=s_0$ ou $s=s_1$,
(ce qui implique $\underline s=\underline s_0$ ou $\underline s=\underline s_1$ donc
$\underline s$ incident à $\underline a$). \deleted{\unsure{Fais|Faux}} Mais si on avait
$s\ne s_0,s_1$, alors $\mathrm{Card}\,a\cap S>2$ (où $S=$
$\mathrm{Im}(\underline S\unsure{-|\to}E)$) a fortiori $\mathrm{Card}\,a\cap C>2$,
où $C$ est la conique circonscrite canonique.
donc $a\subset C$, donc $C$ est \deleted{dégénérée} \inserted{singulière}, donc
$\alpha=-2$. Mais le cas $\alpha=-2$ a été examiné,
et on trouve $C=d\cup d'$, $d$ et $d'$ deux
droites parallèles \uline{distinctes}, avec $s_i\in d$
pour $i$ pair, $s_i\in d'$ pour $i$ impair, donc
$\mathrm{dr}(s_0,s_{\unsure{1|j}})\cap C=\{s_0,s_{\unsure{1|j}}\}$, absurde -- \unsure{gagné}.
% archive sheet number: 105
\scanpage{107}
\uline{Corollaire} Soit $\varphi$ \unsure{repr.|réal.} quelconque de \deleted{\unsure{$\underline G$}}
$\Pi$ polygone combinatoire de degré $N$
($3\leq N\leq+\infty$). \inserted{Soit $\underline u$ un gén de $\mathrm{Aut}^+(\Pi)$, $u$ son image en $\mathrm{Aff}(E)$.} Alors les entiers suivants
sont tous égaux :

\[
\text{1)}\quad\mathrm{Card}\,(S=\mathrm{Im}(\underline S\unsure{-|\to}E))
\]

\[
\text{2)}\quad\mathrm{Card}\,(A=\mathrm{Im}(\underline A\to\mathrm{Dr}(E))
\]

\[
\text{3)}\quad\mathrm{Card}\,(G^+=\mathrm{Im}\,\mathrm{Aut}^{+}(\Pi)\to\mathrm{Aff}(E))
\]

\[
\text{4) Ordre de }u\in\mathrm{Aff}(E)
\]

\deleted{\unsure{Par}} Cet entier \inserted{$n$} est \deleted{$\geq3$ et est} un diviseur de $N$
(condition vide si $N=+\infty$) est $\unsure{\infty|\text{ou}}\geq3$, et c'est
aussi le plus \deleted{grand} \inserted{petit des} diviseurs \inserted{$n$} de $N$ tels
que $\varphi$ se factorise par le \inserted{polygone} quotient
\inserted{$\Pi_{(n)}$} de $\Pi$ d'ordre $n$, et enfin l'unique
diviseur $n$ de $N$ tel que $\varphi$ se factorise
par le polygone quotient $\Pi_{(n)}$ de $\Pi$
d'ordre $n$ \uline{de façon} que la réalisation
obtenue de $\Pi_{(n)}$ soit \uline{fidèle}.

\begin{center}\rule{1.6cm}{0.35pt}\end{center}

Mais il manque un \unsure{sorite|sorte} sur les
morphismes \inserted{(étales)} de polygônes \inserted{comb.} (ou plus
gén\textsuperscript{t} de \unsure{cartes|contours} \inserted{comb.}\ldots), nous permettant
% archive sheet number: 106
\scanpage{108}
de parler du quotient $\Pi_{(n)}$ de $\Pi$,
pour tt diviseur $n$ de l'\unsure{ordre} $N$ de $\Pi$.
% archive sheet number: 107
\scanpage{109}
\margin{(\unsure{18 p|78 p})}
Polynômes cyclotomiques
% archive sheet number: 108
\scanpage{110}
\diagram{\begin{tikzcd}[ampersand replacement=\&,row sep=small]
\mathbb Z[T,T^{-1}] \& {}\\
\mathbb Z(U)\arrow[u,hook] \& U=T+T^{-1}
\end{tikzcd}}

\[
\Sigma_i(U)=T^i+T^{-i}=U^i+\cdots\qquad(i\geq0)
\]
\[
\resizebox{\linewidth}{!}{$\displaystyle\begin{aligned}
F_{2m}(U)&=(1+T^2+\cdots+T^{2(m-1)})T^{-(m-1)}
=\Sigma_{m-1}(U)+\Sigma_{m-3}(U)+\cdots=U^{m-1}+\cdots\quad\big|\,m\geq1\\
F_{2m+1}(U)&=(1+T+T^2+\cdots+T^{2m})T^{-m}
=\Sigma_m(U)+\Sigma_{m-1}(U)+\cdots=U^m+\cdots\quad\big|\,m\geq0
\end{aligned}$}
\]
\[
-\qquad\left\{\begin{aligned}d^\circ(\Sigma_i)&=i\\
d^\circ F_n&=\left[\frac{n-1}{2}\right]=
\begin{cases}m-1&\text{si }n=2m\\m&\text{si }n=2m+1\end{cases}\quad n\geq1
\end{aligned}\right.
\]
\margin{\uline{NB} On pose $\begin{aligned}F_0&=0\\F_n&=F_{-n}\quad\text{si }n<0\end{aligned}$}

\[
\left\{\begin{aligned}
F_{2m}(2)&=m,& F_{2m}(-2)&=(-1)^{m-1}m\\
F_{2m+1}(2)&=2m+1 & F_{2m+1}(-2)&=(-1)^m
\end{aligned}\right.
\]
\[
\left\{\begin{aligned}F_{2m}(U)(T^2-1)&=(T^{2m}-1)T^{-(m-1)}\\
F_{2m+1}(U)(T-1)&=(T^{2m+1}-1)T^{-m}\end{aligned}\right.
\]
\[
F_n(U)F_{n+1}(U)(T-1)(T^2-1)=(T^n-1)(T^{n+1}-1)T^{-(n-1)}
\]
\[
(F_n,F_{n+2}^{\unsure{\prime}})=(1)\qquad n\geq1\quad(\text{et }\unsure{=|\text{aussi}|\text{donc}}\text{ pour }n\in\mathbb Z)
\]
\[
\begin{aligned}
X'_n&=F_n F_{n+1}=U^{n-1}+\cdots\\
Y'_n&=F_n F_{n-1}=U^{n-2}+\cdots
\end{aligned}\quad\bigg|\quad\text{(déf)}
\]
\margin{$\begin{array}{rl}
F_0(U)&=0\\[-2pt]
\multicolumn{2}{l}{\left\{\begin{aligned}F_1(U)&=1\\ F_2(U)&=1\end{aligned}\right.}\\
\multicolumn{2}{l}{\left\{\begin{aligned}F_3(U)&=\deleted{\unsure{1}}U+1\\ F_4(U)&=U\end{aligned}\right.}\\
\multicolumn{2}{l}{\left\{\begin{aligned}F_5(U)&=U^2+U\unsure{-|+}1\\ F_6(U)&=U^2-1\end{aligned}\right.}
\end{array}$}
\margin{$\begin{aligned}
\unsure{-}X'_{-1}&=0\\
X'_0&=0\\ X'_1&=1\\ X'_2&=U+1
\end{aligned}$}

\uline{Prop.} $\forall n\in\mathbb Z$, on a $X'_n-UX'_{n-1}+X_{n-2}=1$

\unsure{NB|DPS} $n\geq1$. Si $n=1$, \unsure{cela|celà} s'écrit $1-U.0+0=1$ OK
Si $n=2$, c'est $U+1-U(1)+0=1$ OK Si $n\geq3$
on écrit
\[
F_n F_{n+1}-U F_{n-1} F_n+F_{n-2} F_{n-1}=1
\]
soit, en multipliant par $(T-1)(T^2-1)T^{n-1}$
\[
\resizebox{\linewidth}{!}{$\displaystyle
(T^n-1)(T^{n+1}-1)-T(T+T^{-1})(T^{n-1}-1)(T^n-1)+T^2(T^{n-2}-1)(T^{n-1}-1)=\deleted{\unsure{1-T^{n+1}-T^n+T^{n-1}}}\inserted{(T-1)(T^2-1)T^{n-1}}$}
\]
\[
\resizebox{\linewidth}{!}{$\displaystyle\begin{aligned}
&\underset{\times}{T^{2n+1}}-\underset{\times}{T^n}-\underset{\times}{T^{n+1}}+\underset{\times}{1}\underbrace{-(T^2+1)(T^{2n-1}-T^n-T^{n-1}+1)}+\underset{\times}{T^{2n-1}}-\underset{\bullet}{T^n}-\underset{\unsure{\bullet|/}}{T^{n+1}}+\underset{\times}{T^2}=\text{\unsure{id}}\\
&\hspace{6em}-\underset{\times}{T^{2n+1}}+\underset{\bullet}{T^{n+2}}+\underset{\times}{T^{n+1}}\unsure{-}\underset{\times}{T^2}-\underset{\times}{T^{2n-1}}+\underset{\times}{T^n}+\underset{\unsure{\bullet|/}}{T^{n-1}}-\underset{\times}{1}
\end{aligned}$}
\]
i.e
\[
T^{n+2}-T^{n+1}-T^n+T^{n-1}=(T-1)(T^2-1)T^{n-1}\qquad\text{OK.}
\]
On en conclut par récurrence sur $n\geq0$
\[
X'_n=X_n\quad\text{i.e.}\qquad X_n=F_n F_{n+1}
\]
puis, comme $X_{-n}=Y_n=\deleted{X_{n-1}}\deleted{\unsure{F_{n-2}}}=F_{n-1}F_n=F_{-n}F_{-n+1}$ \ $(n\geq1)$
% archive sheet number: 109
\scanpage{111}
suffi que le gouvernement amorce une politique hostile aux étrangers\unsure{,|.}
\handwritten{\inserted{Ce sont là les}} boucs émissaires bienvenus (comme jadis les juifs dans l'Allemagne hitlérienne)
\handwritten{\unsure{pour|par}} une récessio\deleted{\unsure{nx}}n économique inexorable, du plus mauvais effet
\deleted{devant l'optimisme de commande officiel et les euphoriques promesses électorales.}
\handwritten{\inserted{face aux euphoriques promesses électorales, alors que le saccage de notre planète continue bon train, ``chez nous'' \deleted{\unsure{tout}} comme \inserted{ailleurs.}}}
Dans l'Allemagne hitlérienne, c'était un crime passible de mort
crématoire que d'être juif, et dangereux pour quiconque de prêter assistance
à un juif. Il en a été de même en France pendant l'Occupation. \deleted{\unsure{La}}
\handwritten{L'}ordonnance\unsure{,} \handwritten{\inserted{du 16.11.1945}}
\deleted{\unsure{loi|lxx} qui aujourd'hui est utilisée comme une arme contre l'étranger (et contre tel citoyen qui \unsure{xdéplaira|déplaira} aux autorités),}
promulguée au lendemain même de l'Occupation, perpé\unsure{tue} bien l'esprit de ces jours maudits\handwritten{! Nous
aurons, nous aussi \illegible{2}\unsure{uifs|juifs}''\ldots}

Combien sont-ils, \unsure{ces nouveaux} Hors-La-Loi du doux pays de France -
les ``étrangers en situation irrégulière'' ? Un million, trois millions ?
Les statistiques officielles les ignorent, par définition - elles font
état seulement des quelque trois millions d'étrangers résidant régulièrement
en France. Si le fameux ``article 21'' est appliqué, ce seront \deleted{\unsure{eux}}\handwritten{là} les
Nouveaux Pariahs de la Société Nouvelle, des hors-la-loi en sursis,
puisqu'il suffira que leur visa arrive à expiration, que leur permis de
séjour ne soit pas renouvellé, pour se trouver eux aussi en ``situation
irrégulière'', au ban de la société. Par un juste retour des choses, ces
trois millions d'étrangers deviendront un véritable danger public - leur
fréquentation peut un beau jour vous valoir la prison\handwritten{!} A quand la ``Solution
finale'' qui nous débarassera de telles sueurs froides ?

\deleted{Si cette ordonnance \unsure{ne vous choque pas, ni ne vous fait} \illegible{1}}
Le signataire de ces lignes a l'honneur, semble-t-il, d'être le premier
``bénéficiaire'' de cette ordonnance passe-partout, se voyant cité à
comparaitre devant le Tribunal Correctionnel de Montpellier, au Palais de
Justice de ladite ville, le 13 Février à 16 heures, poursuivi pour ``Avoir
le 16 Novembre 1975 par aide directe en lui offrant gratuitement nourriture
et logement, facilité le séjour irrégulier \unsure{d'un|dhun} étranger en France, en
l'espèce Kuniomi Masunaga, ressortissant japonais'', `` ces faits constituant
le délit prévu et puni par l'article 21 de l'ord. du 16.11.1945''.

% margin note written sideways (rotated 180 degrees)
\margin{\handwritten{$\begin{aligned}1+T^2+T^4&=(T^{-2}+T^2)+1=U^2-1\\
F_{2m+2}(U)(T^2-1)&=(T^{2\deleted{\unsure{m+1}}m+2}-1)T^{-m}\end{aligned}$}}
% archive sheet number: 110
\scanpage{112}
on trouve
\[
\left\{\begin{aligned}X_n&=F_nF_{n+1}\\Y_n&=X_{n-1}=F_{n-1}F_n\end{aligned}\right.\qquad\bigg|\quad\forall n\in\mathbb Z
\]

Autre démonstration : on voit par induction sur $n\geq1$
que $\deleted{\unsure{F}}X_n(U)$ est \inserted{unitaire} de degré $n-1$ en $U$ (grâce à la relation
$X^{\unsure{\prime}}_n-UX_{n-1}+X_{n-2}=1$ i.e. $X_n=UX_{n-1}-X_{n-2}+1$) \deleted{\unsure{Il a donc}}
tout comme $X'_n(U)$. \deleted{D'ailleurs ils ont les \unsure{mêmes} racines dans $\overline{\mathbb Q}$ -- d'ailleurs celles de $X'_n$ sont les}
\inserted{Considérons les}
\[
\begin{aligned}T^{n-1}X_n(U)&=A_n(T)\\T^{n-1}X'_n(U)&=A'_n(\unsure{T})\end{aligned}
\]
ce sont des polynômes \unsure{unitaires} en $T$ et de $=$
degré $2(n-1)$. les racines du deuxième \inserted{\unsure{dans|de} $\overline{\mathbb Q}$} sont
les \deleted{racines de} $\zeta\in\overline{\mathbb Q}$ telles que $\zeta^n\unsure{=|\in}1$ ou $\zeta^{n+1}=1$,
$\zeta\ne\pm1$ \deleted{\unsure{de}}, en nbre $2n+1-3=2n-2$ -- elles sont
distinctes. On a vu que ce sont les racines
de $A_n(T)$, d'où l'égalité des deux polynômes
($n\geq1$).

\deleted{Description} \inserted{\unsure{Matrice|Matrices}} de $u^n$. Comme $u^n$ transforme
$s_0,s_1,s_{-1}$ en $s_n,s_{n+1},s_{n-1}$ de coordonnées
\[
\resizebox{\linewidth}{!}{$\displaystyle
\deleted{\unsure{F_n F_n}}
\begin{cases}F_n(\alpha)F_{n+1}(\alpha)\\F_{n-1}(\alpha)F_n(\alpha)\end{cases},\
\begin{cases}F_{n+1}(\alpha)F_{n+2}(\alpha)\\F_n(\alpha)F_{n+1}(\alpha)\end{cases},\
\begin{cases}F_{\unsure{n-1|n+1}}(\alpha)F_n(\alpha)\\F_{\unsure{n+2|n-2}}(\alpha)F_{n-1}(\alpha)\end{cases},$}
\]
sa matrice est
\[
\resizebox{.94\linewidth}{!}{$\displaystyle\deleted{u^n=\begin{pmatrix}
F_{n+1}(\alpha)(F_{n+2}(\alpha)-F_n(\alpha))&&F_n(\alpha)F_{n+1}(\alpha)\\
F_n(\alpha)(F_{n+1}(\alpha)-F_{n-1}(\alpha))&&F_{n-1}(\alpha)F_n(\alpha)\\
0&0&1\end{pmatrix}}$}
\]
\[
\resizebox{\linewidth}{!}{$\displaystyle u^n=\begin{pmatrix}
F_{n+1}(\alpha)(F_{n+2}(\alpha)-F_n(\alpha))&F_n(\alpha)(F_{n-1}(\alpha)-F_{n+1}(\alpha))&F_{\unsure{n|n+1}}(\alpha)F_{\unsure{n+1}}(\alpha)\\
F_n(\alpha)(F_{n+1}(\alpha)-F_{n-1}(\alpha))&F_{n-1}(\alpha)(F_{n-2}(\alpha)-F_n(\alpha))&F_{n-1}(\alpha)F_n(\alpha)\\
0&0&1
\end{pmatrix}$}
\]
% archive sheet number: 111
\scanpage{113}
\noindent\(\left\{\vcenter{\hbox{\begin{minipage}{.91\linewidth}
\uline{Théorème} Soit \deleted{\unsure{un}} $k$ anneau, $\alpha\in k$, d'où $u_\alpha:k^{2}\to k^2$
\deleted{transformation} \inserted{\unsure{autom.|automorph.}} affine
\[
u_\alpha(x,y)=(1-y+\alpha x,x)
\]
Satisfaisant, \deleted{\unsure{si}} posant $s_n=u_\alpha^n(s_0)$ (où $s_0=(0,0)$)
\[
\begin{aligned}
\deleted{\unsure{u_\alpha^n(s_0)|u_\alpha(s_0)}}\qquad s_1&=u_\alpha^{\deleted{\unsure{n}}}(s_0)=(1,0)\\
s_{-1}&=u_\alpha^{-1}(s_0)=(0,1)\quad\text{i.e.}\quad u_\alpha(0,1)=(0,0).
\end{aligned}
\]
Soit $n\in\mathbb Z$. Conditions équivalentes

a) $u_\alpha^n=\mathrm{id}$

b) $s_n=s_0$ i.e $u_\alpha^ns_0=s_0$

b$_{\unsure{i|i+}}$) (si $i\in\mathbb Z$ fixé) $u_\alpha^ns_i=s_i$

d) $F_n(\alpha)=0$
\end{minipage}}}\right.\)

a $\Rightarrow$ b trivial

b $\Rightarrow$ a car b) implique \inserted{(en multipliant par $u_\alpha^i$)} $u_\alpha^ns_i=s_i$ $\forall i$, en
partic si $i=\pm1$ on trouve $\begin{cases}u_\alpha^n(s_{-1})=s_{-1}\\u_\alpha^n(s_1)=s_1\end{cases}$
d'où $u_\alpha^n=\mathrm{id}$
\[
\begin{aligned}
\text{b)}&\Rightarrow\text{b}_i\text{)}\quad\text{en multipliant par }u_\alpha^i\\
\text{b}_i\text{)}&\Rightarrow\text{b)}\qquad\rule{2.8cm}{.3pt}\qquad u_\alpha^{-i}
\end{aligned}
\]
donc $a\Longleftrightarrow b\Longleftrightarrow c$

Comme $s_n=(F_n(\alpha)F_{n+1}(\alpha),F_{n-1}(\alpha)F_n(\alpha))$, on voit
que d) $\Rightarrow$ b) i.e. $F_n(\alpha)=0\Rightarrow s_n=s_0\ (=(0,0))$
Inversement, si $F_n(\alpha)F_{n+1}(\alpha)=0$, $F_n(\alpha)F_{n-1}(\alpha)=0$ 
il s'ensuit $F_n(\alpha)(a F_{n-1}(\alpha)+b F_{n+1}(\alpha))=0$ $\forall a,b\in k$.
Or $\exists A_n(U),B_n(U)\in\mathbb Z[U]$ tels que
\[
A_n(U)F_{n-1}(U)+B_n(U)F_{n+1}(U)=1,
\]
d'où, faisant $a=A_n(\alpha)$, $b=B_n(\alpha)$, la relation $F_n(\alpha)=1$,
cqfd\ldots
% archive sheet number: 112
\scanpage{114}
\uline{NB} On peut prendre
\[
\left\{\begin{aligned}
\deleted{\unsure{A}}B_n(U)&=F_n(U)\\
\deleted{\unsure{B}}A_n(U)&=F_{n-2}(U)-UF_{n-1}(U)
\end{aligned}\right.
\]
% archive sheet number: 113
\scanpage{115}

\[
\deleted{T^n-1}
\]
\[
\Phi_n(T)=\prod_{\substack{\xi\ \text{d'ordre }n\\ \text{dans }\overline{\mathbb Q}}}(T-\xi)
\quad\text{polynôme \inserted{$\in\mathbb Z[T]$} unitaire de degré }\varphi(n)=\operatorname{card}(\mathbb Z/n\mathbb Z)^*
\]
\[
\begin{aligned}
\varphi(\underbrace{p_1^{\alpha_1}\cdots p_r^{\alpha_r}}_{n})
 &=\prod_i\varphi(p_i^{\alpha_i})=\prod p_i^{\alpha_i}-p_i^{\alpha_i-1}\\
 &=n\prod\left(1-\frac1{p_i}\right)
\end{aligned}
\]
\[
T^n-1=\prod_{\substack{d\in\mathbb N^*\\ d\mid n}}\Phi_d
\]
\[
\text{donc}\qquad
\Phi_n=(T^n-1)\big/\prod_{\substack{d\in\mathbb N^{\times}\\ d\mid n,\,d\ne n}}\Phi_d
\]

Si $k$ est un corps de car. \uuline{$p\nmid n$} (donc les racines de
$T^n-1=0$ y sont distinctes\ldots) alors les racines de
$\Phi_n(T)=0$ dans $k$ sont les racines \uline{primitives} $n$.ièmes
de $1$ dans $k$.

Ainsi si \inserted{\unsure{c|c'|on a}} \deleted{\unsure{on}} une réalis. \unsure{lin.}
géom. \inserted{rég.} de $\Pi_\infty$ sur $k$
d'inv. $\alpha$ tel que $\alpha=\zeta+\zeta^{-1}$,
pour qu'elle se factorise en une réalisation géom.
\uline{fidèle} de $\Pi_n$ sur $k$, \inserted{lorsque $n\geq3$,} il f. et s. que on ait
$\Phi_n(\zeta)=0$ i.e. $\Psi_n(\alpha)=0$, où $\Psi_n\in\mathbb Z[U]$ est
donné par
\[
\Psi_n(T+T^{-1})=T^{-\varphi'(n)}\Phi_n(T)
\qquad\substack{\varphi'(n)=\tfrac12\varphi(n)\\ \text{si }n\geq 3}
\]

Donc si $n$ \inserted{$\geq3$} n'est ni premier, ni double d'un \deleted{\unsure{premier}}
nb premier (donc $n=8,9,12,15,16\ldots$) alors l'anneau
modulaire pour les réalisations géom. \uline{fidèle}
\margin{$\frac{\varphi(n)}2$}
% The black disc in the interword space has no decipherable character.
% Three lines in the opposite orientation, cancelled by four strokes.
\par\noindent\rotatebox{180}{\begin{minipage}{.82\linewidth}
\deleted{Liebe Dagmar,}\\
\deleted{Vielen Dank für Deine}\\
\deleted{lieben Geburtstagsgrüsse.}
\end{minipage}}
% archive sheet number: 114
\scanpage{116}

de $\deleted{\unsure{\Phi}}\Pi_n$ (pol. comb. à $n$ côtés), \deleted{\unsure{au moins}} est
\[
k'_n=\mathbb Z\bigl[\tfrac1n\bigr][U]/\Psi_n[U]
\]
$\alpha'_n\in k'_n$ étant l'image canonique de $U$.

Si $n=\deleted{\unsure{2p}}p$ ($p$ premier\inserted{ impair}\deleted{\unsure{impair}}) \deleted{$p$ et $k$ corps de car. $p$,
alors l'invariant du polygone régulier à $2p$ côtés sur $k$
est $+2$, le polygone cyclotomique}
\deleted{\inserted{D'ailleurs \unsure{si} $p$ impair}}
\[
\deleted{\Phi_{2p}(T)=\Phi_p(-T)\deleted{\unsure{{}^{p}}}
 =\deleted{\unsure{(T^{p-1})(T+1)}}T^{p-1}-T^{p-2}\ldots+1\deleted{\unsure{(T+1)^{p-1}}}}
\]
\deleted{\deleted{\unsure{En car. $p$ on}} En car. $p$, ses solutions sont \deleted{des} toutes
\deleted{$\Psi_{2p}(U)=$} confondues, et égales : $-1$ (celles de $\Phi_p$ étant égales
: $1$) donc les solutions de $\Psi_{\deleted{\unsure{p}}2p}$ sont égales : $1$,
les racines de $\Phi_p(T)=T^{p-1}\ldots1$ \deleted{\unsure{donc alors toutes égales : $1$}}
donc les racines de $\Psi_p(T)$ sont \unsure{toutes} égales : $2$,
Donc si $\alpha\in k$, la représentation de \unsure{donc|dont} celles de
$\Psi_p(U)$ ($=F_p(U)$ ici !) égales : $2$.}
toute réalisation géom. régulière de $\Pi_p$ sur un corps est fidèle
--- d'\deleted{fait}ailleurs on trouve
\[
\begin{aligned}
\Phi_p(T)&=\frac{T^p-1}{T-1}=T^{p-1}\ldots+1\\
\Psi_p(\deleted{\unsure{T}}U)&=F_p(\deleted{\unsure{T}}U)
\end{aligned}
\]
et on trouve l'anneau modulaire
\[
\underline{k}'_p=\mathbb Z[U]/\Psi_p(U)=\mathbb Z[U]/F_p(U)
\]

Supposons enfin que $n=2p$ ($p$ premier). Si $p$
impair, \deleted{\unsure{on a facilement que} ($\zeta\in\overline{\mathbb Q}$) : \unsure{le fait|à la fois}
$(\deleted{\unsure{\ne}}\zeta^2\unsure{-1|\ne1})^p=1$
donc les racines de $\Phi_{2p}(T)=0$ \unsure{sont} les $\zeta$ tels que $(\zeta^2)^p=1$
$\zeta^2\ne1$, $\zeta^p\ne1$ sont aussi les $\zeta\in\overline{\mathbb Q}$ tels que
$(\zeta^2)^p=1$, $\zeta^2\ne1$}
\[
\deleted{\begin{aligned}
 & (T^{2p}-1)\big/\unsure{(T^2-1)(T-1)|(T-1)(T^p-1)}(T^{p-1}\ldots+1)\\
 & \zeta^{2p}=1,\ \zeta^{\unsure{1}}\ne1,\ \zeta^p\ne1,\
 \unsure{(\zeta^2)^p\ne1|\zeta^p\ne1},\ \zeta^2\ne1,\ (T\unsure{+1}),\ \zeta^{2p},\ \zeta^p\ne1
\end{aligned}}
\]
% archive sheet number: 115
